% !TeX root = ../../../main.tex

\subsection{Selected Error Analyses}\label{ch05:apd:selected-errors}

\ChFiveHeadingOnly

\subsubsection{Inter-LLM Comparison of Temporal Annotations}\label{ch05:apd:inter-llm-comparison}

To illustrate qualitative differences among annotators, we compare temporal event sequences generated by several LLMs for a representative case report (PMC10629858) \citep{abu2023rare}, which describes a patient with lepromatous leprosy who later developed abdominal complications, recurrent bacteremia, severe sepsis, and ultimately died \citep[Supplementary Material: Inter-LLM Comparison of Temporal Annotations]{noroozizadeh-pmoa-tts-v4}.

Across models, the principal clinical trajectory is consistently recovered: confirmation of lepromatous leprosy, antibiotic treatment, evolving abdominal symptoms with corresponding imaging findings, subsequent sepsis, ICU readmission, and death.
All annotators preserve the
temporal polarity of events, assigning negative times to historical information and positive times to post-admission findings.

Differences emerge in event granularity and phrasing.
Some models split conjunctive phrases into multiple events (e.g., listing each component medication separately), whereas others
compress them into single entries.
Several annotators also produce near-duplicate events at the same timestamp (e.g., both “diagnosed with lepromatous leprosy’’ and “skin biopsy confirmed’’), reflecting variation in how models segment clinically related information within the narrative.

More pronounced differences appear in temporal normalization.
For the expression “two months before admission,” models produce hour offsets ranging from approximately $-1464$ to $-1344$, and one model assigns $-144$ due to an apparent day–hour unit error.
Long-range events such as the autologous bone marrow transplant plan, ICU readmission, and death also receive timestamps that vary by hundreds of hours across annotators.
Despite this variability in absolute timing,
the relative ordering of major clinical milestones remains largely preserved.

A recurring issue across models involves distinguishing planned actions from events that actually occurred.
For example, most annotators correctly represent “planned autologous bone-marrow transplant’’ with an associated future timestamp, which is appropriate because planned interventions can be temporally located and may influence subsequent decision-making.
However, several models omit the word “planned,’’ converting the entry into “autologous
bone-marrow transplant’’ and thereby representing an intended but unperformed procedure as an event that occurred.
In such cases, the discrepancy arises from the event label rather than the timestamp itself, illustrating how small linguistic omissions can alter the extracted temporal semantics.

% !TeX root = ../../../main.tex
\begin{table}[htbp]
\caption{Common error types observed across LLM-generated temporal annotations.}\ChFiveFloatLabel{ch05:tab:error_taxonomy}
\ChFiveTableBodyStart
\begin{ChFiveFloatBody}

\centering


\begin{tabular}{>{\raggedright\arraybackslash}p{3.2cm} >{\raggedright\arraybackslash}p{9cm}}
\toprule
\textbf{Error Type} & \textbf{Description and Examples} \\
\midrule

\textbf{Event granularity mismatch} 
& Models differ in splitting vs.\ compressing multi-component statements. 
For example, medications may appear as three separate events (``rifampicin'', ``clofazimine'', ``dapsone'') or as one combined entry. \\[6pt]

\textbf{Near-duplicate events} 
& Multiple paraphrased versions of the same clinical concept assigned to the same timestamp, such as 
``diagnosed with lepromatous leprosy’’ and 
``skin biopsy confirmed lepromatous leprosy.’’ \\[6pt]

\textbf{Temporal normalization variability} 
& Large differences in converting narrative expressions (e.g., ``two months before admission’’) to numerical timestamps, with offsets ranging from $-1464$ to $-1344$ hours and occasional scaling errors (e.g., $-144$ hours). Long-range events (e.g., transplant plan, ICU readmission, death) also vary by hundreds of hours across models. \\[6pt]

\textbf{Intent vs.\ occurrence confusion} 
& Some models omit modality markers such as ``planned,'' converting intended future events into events marked as completed. 
The timestamp itself may be reasonable; the error lies in the event label. \\[6pt]

\textbf{Phrasing and specificity variation} 
& Models differ in the level of detail used in event descriptions.  
Some outputs retain clinically precise phrases (e.g., ``mural thickening of the terminal ileum’’), while others produce generic events (e.g., ``abdominal abnormality’’). \\[6pt]

\textbf{Long-range timestamp degradation} 
& Accuracy decreases for events far from the admission anchor, with larger dispersion in timestamp predictions for events weeks or months after presentation. \\
\bottomrule
\end{tabular}


\end{ChFiveFloatBody}
\end{table}


Overall, the models capture the broad sequence of events but differ in how they segment conjunctive statements, handle near-duplicates, and resolve absolute times—particularly for
long-range events.
A taxonomy of common annotation errors observed across models—including granularity differences, duplicate events, temporal normalization variability, and intent–occurrence misclassification—is summarized in Table~\ref{ch05:tab:error_taxonomy}.

% !TeX root = ../../../main.tex

\subsubsection{Event Counts and Annotation Schema}\label{ch06:apd:error-analysis}
To investigate the discrepancies in event match recall between the i2m4 and \texttt{sepsis-10} datasets, we analyzed the datasets' structure and examined model performance using Llama 3.3 as a representative model, as the pattern of differences observed across datasets was consistent for all models.

\ChSixParagraph{Overview of Clinical Event Identification.}
For the i2m4 dataset, manual annotations by clinicians identified an average of $195\pm94$ clinical events per report across $20$ case reports, compared to $80\pm28$ clinical events per report for \texttt{sepsis-10} across $10$ case reports.
Llama 3.3 identified an average of $120\pm61$ clinical events for i2m4 and $78\pm24$ clinical events for \texttt{sepsis-10}.
These results highlight a systematic bias of under-identification of clinical events in i2m4 compared to \texttt{sepsis-10} by Llama 3.3.

Additionally, the i2m4 dataset contains approximately $2.4\times$ more clinical events per report on average than \texttt{sepsis-10}, as annotated by the clinician, with a much larger variance in event counts across reports.
This heterogeneity reflects structural differences in the datasets: i2m4 reports often include a large number of structured entries, whereas \texttt{sepsis-10} reports are only made-up of free text.
Llama 3.3 over-identified clinical events (i.e., predicted more events than the manual annotators) in $15\%$ of i2m4 reports, while under-identifying events in $85\%$ of i2m4 reports.
For \texttt{sepsis-10}, the model showed a more balanced distribution, over-identifying and under-identifying events in $50\%$ of reports each (Table~\ref{ch06:tab:event_identification_bias}).
These patterns suggest that the higher density and variability of clinical events in i2m4 pose a greater challenge for Llama 3.3.






% !TeX root = ../../../main.tex
\begin{table}[htbp]
\caption{Proportion of reports where Llama 3.3 over- or under-identified clinical events compared to manual annotations. i2m4 shows a strong bias toward under-identification, whereas \texttt{sepsis-10} has a more balanced distribution.}\ChSixFloatLabel{ch06:tab:event_identification_bias}
\ChSixTableBodyStart
\begin{ChSixFloatBody}

\setlength{\tabcolsep}{4pt}
\centering
\begin{tabular}{|l|c|c|}
\hline
\textbf{Dataset} & \makecell{\textbf{\% Reports} \\ \textbf{Over-Identified}} & \makecell{\textbf{\% Reports} \\ \textbf{Under-Identified}} \\ \hline
\textbf{i2m4}    & 15\%                                & 85\%                                 \\ \hline
\textbf{\texttt{sepsis-10}} & 50\%                              & 50\%                                 \\ \hline
\end{tabular}



\end{ChSixFloatBody}
\end{table}


\ChSixParagraph{Dataset Structure Differences.}
The differences in dataset structure between i2m4 and \texttt{sepsis-10} are also evident in the length of clinical events.
As shown in Table~\ref{ch06:tab:character_length}, the manually annotated clinical events in i2m4 are significantly shorter, averaging $10\pm7$ characters, compared to $28\pm16$ characters in \texttt{sepsis-10}.
This reflects the structured nature of i2m4 reports, which often include tables with abbreviations (e.g., ``K'' for Potassium, ``Na'' for Sodium).



% !TeX root = ../../../main.tex
\begin{table}[htbp]
\caption{Average length of clinical events in characters for both datasets. Manual annotations in i2m4 are significantly shorter due to frequent abbreviations and structured text from MIMIC tables, while Llama 3.3 generates events of comparable length across datasets.}\ChSixFloatLabel{ch06:tab:character_length}
\ChSixTableBodyStart
\begin{ChSixFloatBody}

\centering
\begin{tabular}{|l|c|c|}
\hline
\textbf{Dataset} & \textbf{Manual Annotator} & \textbf{Llama 3.3} \\ \hline
\textbf{i2m4}    & $10\pm7 $                              & $17\pm14$                        \\ \hline
\textbf{\texttt{sepsis-10}} & $28\pm16$                            & $21\pm13$                        \\ \hline
\end{tabular}



\end{ChSixFloatBody}
\end{table}


Llama 3.3-generated clinical events are of similar length across datasets, with averages of $17\pm14$ characters for i2m4 and $21\pm13$ characters for \texttt{sepsis-10}.
This consistency suggests that Llama 3.3 applies similar generative behavior across datasets, often producing expanded forms of abbreviations in i2m4.
While this generative behavior mitigates variability in Llama 3.3 outputs, it creates a mismatch with the shorter, more concise manual annotations in i2m4.
This mismatch can then contribute to higher cosine distances between embeddings, lowering the event match recall for i2m4.


The complete annotation galleries and additional event-category, adjusted-denominator, and outlier-filtered analyses remain in the selected CHIL manuscript, Appendices H and L \citep{pmlr-v333-noroozizadeh26a}.
